\section{Plano de Análise}

\textbf{5.1 Limpeza} (nesta ordem; contagem de descartes reportada em fluxo brutos $\to$ válidos): (1) ITEM\_18 $\neq 2$; (2) tempo total $< 120\text{ s}$ (limiar fixo, $\sim 5\text{ s/item}$, quando mensurável por complemento de temporização; caso o Google Forms nativo registre apenas o carimbo de envio sem log de início, o descarte por falta de atenção baseia-se nos critérios 1 e 3, reportando-se a limitação da plataforma); (3) desvio-padrão zero entre os ITEM\_08--17 (resposta uniforme / \textit{straightlining}); (4) vetores de resposta completos idênticos (duplicatas; sem login não é possível impedi-las, limitação declarada); (5) \enquote{Não se aplica} = faltante: o escore do construto é a média dos itens respondidos se $\ge 75\%$ deles foram respondidos, senão exclusão \textit{pairwise}, sem imputação; (6) inversão do ITEM\_16 ($6 - x$).

\textbf{5.2 Confiabilidade.} Alfa de Cronbach (IC 95\% por \textit{bootstrap}) para C\_VEL e C\_QUAL; correlação de Spearman ($\rho$) entre ITEM\_19 e ITEM\_20 para M\_REV. Aceitação: $\alpha \ge 0{,}70$ para C\_VEL e C\_QUAL, e $\rho \ge 0{,}50$ para M\_REV. Se $\alpha < 0{,}70$, remove-se no máximo um item, apenas se sua correlação item-total for $< 0{,}30$ e o $\alpha$ sem ele subir $\ge 0{,}05$; caso contrário (ou se $\rho < 0{,}50$ em M\_REV), QP3 e QP4 são analisadas item a item (Bonferroni) e a limitação é relatada. Com $n \approx 60$ não há análise fatorial.

\textbf{5.3 Testes.}
\begin{itemize}
  \item \textbf{QP1 e QP2:} mediana, moda, desvio absoluto mediano (MAD) e \% de concordância (respostas 4--5) com IC de Wilson de 95\%, por item e construto; gráficos de barras divergentes.
  \item \textbf{QP3:} Kruskal-Wallis, separado para C\_VEL e C\_QUAL, com $\alpha = 0{,}025$ cada (Bonferroni, família de 2). Se $H_0$ for rejeitada: pós-teste de Dunn com Bonferroni para 3 pares (Jr--Pl, Jr--Sr, Pl--Sr). Tamanho de efeito: $\varepsilon^2$ (épsilon ao quadrado) com IC por \textit{bootstrap}. Se algum nível tiver $< 8$ respondentes, fundem-se os adjacentes menores (e.g., Jr+Pl vs.\ Sr) e aplica-se Mann-Whitney (efeito $r$), registrando a fusão.
  \item \textbf{QP4:} Spearman ($\rho$) entre M\_REV e C\_QUAL, bicaudal, $\alpha = 0{,}05$, IC por \textit{bootstrap}; efeito segundo Cohen (0,10 fraco; 0,30 moderado; 0,50 forte). Sensibilidade (exploratória): correlação parcial controlando M\_SEN e M\_USO; C\_QUAL sem o ITEM\_16 e, separadamente, sem o ITEM\_17. Correlação não implica causalidade.
  \item \textbf{Exploratórias:} distribuição de M\_TOOL e M\_USO; associação ferramenta $\times$ senioridade por $\chi^2$ (ou teste exato de Fisher se houver células esperadas $< 5$).
\end{itemize}

\textbf{5.4 Respostas abertas (ITEM\_21--22).} Análise temática indutiva por dois pesquisadores independentes (riscos, e.g., débito técnico oculto, dependência cognitiva; benefícios, e.g., prototipagem rápida, menos \textit{boilerplate}). Ao menos 20\% das respostas são codificadas por ambos, com $\kappa$ (Kappa) de Cohen (meta $\ge 0{,}60$; abaixo disso, refina-se o livro de códigos); divergências são resolvidas por consenso e registradas. Trechos com nomes de pessoas, empresas ou projetos são suprimidos antes de qualquer divulgação.
